% Chapter 6: Experimental Evaluation
\section{Evaluation Objectives}

The evaluation aimed to assess both the technical effectiveness of the framework and its pedagogical impact on student learning.

\section{Research Methodology}

\subsection{Study Design}

\OpenQuestionPrompt
  {Evaluation strategy}
  {What evaluation design best fits the current maturity of the platform: expert heuristic review, pilot classroom deployment, comparative usability study, design-based research iteration, or a staged mixed-methods sequence? State why that design is appropriate for a knowledge API with both course-oriented and gamified clients.}

\subsection{Participants}

\CategoryPrompt
  {Sampling and recruitment}
  {Describe who should participate in the study and how they should be grouped. Make explicit whether the participants are middle-school learners, teachers, curriculum experts, university pilot users, or a mixture of these audiences.}
  {learners by age band; teachers by subject background; expert reviewers; prior CS exposure; assignment to conditions}

\section{Metrics and Measurements}

\subsection{Learning Outcome Metrics}

- Pre/post assessment scores measuring knowledge acquisition
- Performance on conceptual understanding tasks
- Transfer of knowledge to novel problems

\subsection{Engagement Metrics}

- Time spent in learning activities
- Frequency and depth of system interactions
- System feature utilization rates

\subsection{Usability Metrics}

- System Usability Scale (SUS) scores
- Task completion rates
- Error rates and recovery patterns

\begin{table}[H]
    \centering
    \begin{threeparttable}
        \caption{Evaluation reporting scaffold for the dissertation prototype.}
        \label{tab:evaluation-reporting-scaffold}
        \begin{tabularx}{\linewidth}{>{\raggedright\arraybackslash}p{0.18\linewidth} Y Y Y}
            \toprule
            Dimension & Example instrument or artifact & Data source & Recommended analysis \\
            \midrule
            Learning outcomes & Pre/post concept checks, guided tasks, short quizzes & Assessment submissions and rubric scores & Gain scores, item-level analysis, and concept-specific interpretation \\
            Usability & \ac{SUS}, task completion, error logs & Surveys, observation notes, screen recordings & Descriptive statistics with issue prioritization \\
            Engagement & Session counts, dwell time, mission completion, revisit rate & Interaction logs from web and mobile clients & Trend summaries and behavioral comparison across services \\
            Pedagogical alignment & Teacher review rubric, standards crosswalk, lesson-plan audit & Expert review sheets and curriculum mappings & Thematic synthesis with alignment gaps highlighted \\
            Trust and transparency & Explanations of recommendations and navigation confidence & Interviews, think-aloud sessions, open-ended survey items & Qualitative coding and representative quotations \\
            \bottomrule
        \end{tabularx}
    \end{threeparttable}
\end{table}

\section{Results}

\subsection{Quantitative Findings}

\EvidencePrompt
  {Minimum quantitative reporting}
  {Report descriptive statistics for each major metric, state the comparison or baseline being used, and record whether the web app, mobile app, and landing-page journeys are being evaluated together or separately. Include effect sizes or confidence intervals whenever the study design supports them.}

\subsection{Qualitative Findings}

\CategoryPrompt
  {Qualitative coding frame}
  {Summarize the recurring themes that emerge from interviews, think-aloud sessions, and teacher feedback. The categories should remain open enough to change as the product evolves, but specific enough to guide consistent coding.}
  {orientation and navigation; motivation and engagement; clarity of lesson plans; trust in recommendations; perceived learning value; client-facing clarity}

\section{Discussion}

\subsection{Interpretation}

The evaluation results demonstrate both strengths and areas for improvement in the framework design and implementation.

\subsection{Limitations}

\OpenQuestionPrompt
  {Scope and generalization}
  {Which limitations are most important to acknowledge at the current stage: sample size, participant population, prototype fidelity, incomplete content coverage, classroom duration, or the difference between a course-oriented web client and a gamified mobile client?}
